\scanpage{100}
Conditions nécessaires pour qu'un $\lambda\in\bar Q$ soit valeur
propre de \deleted{\unsure{son|un}} frobenius, pour une variété projective
lisse définie sur le corps $\mathbb F_q$, $q=p^n$ :

\deleted{\[\left\{\begin{aligned}
&1)\ \lambda\in Q'(i)\\
&2)\ \deleted{\unsure{\lambda}}\lambda^2\ \text{est de la forme }q^i u,\quad\text{où }u\in Q'(i)\text{ est une unité \uline{exc. $p$} ;}\\
&\text{Si }p,p'\text{ sont deux idéaux premiers de}
\end{aligned}\right.\]}

\[
\begin{aligned}
1)&\quad\boxed{\lambda\ \text{entier algébrique}}\\
2)&\quad\boxed{\lambda\in Q'(\sqrt{-1})}\\
3)&\quad\lambda^{\unsure{1|{}|2}}=q^{i/2}u,\quad\text{avec \inserted{$i\in\mathbb Z$,} $u$ une unité \inserted{absolue} exc. $p$}
\end{aligned}
\]
(Cela détermine évidemment de façon unique $i$ et $u$)
\margin{N.B. $q^{i/2}\in Q'$ !}
\[
\Updownarrow\quad \boxed{\exists\ i\ \text{tel que }|\lambda_\alpha|=q^{i/2}\text{ pour }\forall\text{ conjugué \unsure{$\alpha$|$\lambda_\alpha$} de }\lambda}
\]
\deleted{\illegible{2}}
\margin{\uline{N.B.} Les conditions 2) 3) 4) sont ``absolues'', la condition 1) est seulement \unsure{diophantienne|restrictive} \ldots\ On désigne par $M(q^{\unsure{n|{}}})$ le ss-groupe de $Q'(i)^*$ défini par \deleted{\unsure{2)|3)}} 3) 4).}

\deleted{4) Si $\mathfrak p,\mathfrak p'$ sont deux idéaux premiers de $Q'(\sqrt{-1})$
induisant le \unsure{m} idéal premier de $Q'$, on a
\[
v_{\mathfrak p}(u)+v_{\mathfrak p'}(u)=0\qquad \text{i.e.}\quad
\boxed{v_{\mathfrak p}(\lambda)+v_{\mathfrak p'}(\lambda^{\unsure{*|{}}})=\frac{in}{2}}
\]
$v_p$ normalisé tel que $v_p(p)=1$.}
\margin{\begin{tikzpicture}[x=1cm,y=1cm]
\node[ellipse,draw,inner sep=1pt,align=center,font=\scriptsize] (c) at (0,0) {Conséquence\\de \uline{2 et 3}\\!!!};
\draw[->,double] (c.north east) -- ++(.72,.42);
\end{tikzpicture}}
% circled, with an arrow to the boxed valuation formula
L'entier $i$ qui intervient ici est appelé le \uline{poids} de $\lambda$
(relativement à $q=p^n$). \deleted{\unsure{Le fait que le poids} \illegible{2}}
\deleted{\unsure{est de}} \inserted{Si les conditions} envisagées sont satisfaites \deleted{\unsure{sur}} pour
$q'=q^m$, elles le sont pour $q$, et le poids $i$
rel. à $q$ est égal à $mi'$, où $i'$ est le poids par
rapport à $q'$.

\margin{Th. d'existence de motifs sur $\mathbb F_q$}
% slanted
\uline{Conjecture} \deleted{\unsure{ou plutôt}} \uline{\unsure{forte|fait|fondamentale}} \uline{les conditions nécessaires
précédentes sont suffisantes}

Cela déterminerait donc entièrement, en termes
arithmétiques simples, le progroupe de Galois \inserted{$G$} motivique
des $\mathbb F_q$. Vérifions que cela donne bien
% archive sheet number: 99

